\chapter{Introduction}\label{ch:intro}

\section{Research background}

A domestic intrusion alarm has one obligation. When a sensor detects an intrusion, a siren must
sound within a bounded time. Everything else the device does — reporting state to a phone,
accepting remote commands, logging events — is secondary, and none of it excuses a failure to
make a noise.

The hardware on which these devices are now built makes that obligation harder to meet than it
first appears. A decade ago an alarm panel was a microcontroller with a dedicated radio module
or a wired connection. Today the entire device is typically one inexpensive
system-on-chip. On the Espressif ESP32, which is the platform studied here, a single die carries
two processor cores, the 802.11 radio and its closed-source driver, the TCP/IP stack, the TLS
implementation, the MQTT client and the alarm application itself. They share the scheduler, the
memory bus and the interrupt controller.

This matters because it creates a path from the outside world to the safety function that does
not exist in a system with separated hardware. An attacker who cannot reach the alarm's logic
directly may still be able to load the shared resources enough to delay it. The attacker needs
no credentials, no access to the building and no knowledge of the application: radio-layer
frames are processed beneath the IP stack whether or not any socket is listening.

\subsection{Embedded safety firmware}

The term \emph{safety firmware} is used in this report for the portion of an embedded program
whose correct and timely execution is the reason the device exists. In an alarm this is the
sensor-conditioning code, the state machine that decides whether the system is armed, and the
output driver that energises the siren.

Two properties distinguish safety firmware from the rest of a program. First, its obligation is
a \emph{bound}, not an average: an alarm that usually sounds in 150\,ms and occasionally takes
four seconds has failed, even though its mean is excellent. Second, it must continue to work
when the rest of the system does not. A loss of network connectivity, a broker outage or a
saturated radio must not prevent the siren from sounding, because the circumstances under which
an alarm is most needed are precisely the circumstances in which the network is least reliable.

Developers of such firmware on dual-core microcontrollers are given two standard pieces of
advice, which this work takes as its subject:

\begin{enumerate}[leftmargin=1.6em,itemsep=4pt]
  \item \textbf{Pin the safety task to its own core}, and place everything network-facing on
        the other, so that network work cannot delay safety work.
  \item \textbf{Gate risky features behind build-time configuration}, so that a dangerous
        capability is compiled out unless the build deliberately asks for it.
\end{enumerate}

Both are reasonable. Both are also \emph{assertions} about the system rather than
demonstrated properties — the first asserts something about timing, the second something about
security — and in ordinary practice neither is tested. That gap is the subject of this work.

\subsection{MQTT-based communication}

MQTT is the dominant application protocol for this class of device. It is a publish--subscribe
protocol in which clients connect to a central broker, publish messages to topics, and
subscribe to topics they wish to receive. Its appeal for constrained devices is real: the
protocol is small, the client libraries are small, and a device behind a domestic router needs
no inbound connectivity, because it establishes an outbound connection to the broker and
receives commands over it.

The architectural consequence is less often discussed, and it is central to
Chapter~\ref{ch:security}. The broker is not a transparent pipe. It terminates each client's
transport connection and \emph{re-originates} every message it forwards. A device receiving a
command therefore receives it from the broker, over the device's own TLS connection to the
broker, with no cryptographic evidence of which client originally published it. Transport
security protects each hop; it conveys nothing about provenance across the relay. MQTT 3.1.1
compounds this by providing no return code meaning ``refused by policy'': a broker that
discards a publication for violating an access control list still acknowledges it, so the
publisher cannot distinguish success from silent refusal.

\subsection{Security challenges in connected IoT systems}

Three difficulties recur in this device class and all three appear in the results reported
here.

\textbf{Credentials are shared with the device.} A device that authenticates to a broker must
hold a credential, and that credential lives in flash on a physically accessible object. The
same is true of any application-layer key. If both are in the same image, an adversary who
extracts one extracts the other — a point that turns out to determine which of two defences
survives which compromise (Section~\ref{sec:sec-flashdump}).

\textbf{The controlling host is part of the trust boundary.} A phone or home-automation server
that legitimately issues commands must hold whatever authority is needed to issue them. Its
compromise therefore yields that authority, and no amount of care inside the device can prevent
it. Chapter~\ref{ch:security} reports this as the configuration that defeats the strongest
design evaluated here.

\textbf{Availability and integrity are separate properties and trade against each other.} A
rate limiter improves resistance to flooding and degrades legitimate availability. A static
shared token preserves availability and provides no replay resistance. The ablation in
Section~\ref{sec:res-ablation} measures this independence rather than assuming it.

\subsection{Why reliable intrusion alarms matter}

The case for taking this seriously does not rest on sophistication. These devices are cheap,
numerous, installed by non-specialists and rarely updated. A homeowner who installs one has
made a decision about physical security on the strength of an implicit promise about timing and
authority. An alarm that can be silenced from outside the building by an adversary with a
leaked credential, or delayed past the point of usefulness by radio-layer noise, has broken
that promise in a way that is invisible until it matters.

The installed base is also large enough that a defect in the standard engineering advice is
worth more attention than a defect in any single product.

\section{Motivation}

The direct motivation is a measurement by Behnke et al.\ \cite{behnke2020}, who instrumented a
periodic critical task on an ESP32 and subjected the device to IP packet load. They found
measurable degradation: the task lost between 2.15\% and 3.8\% of its cycles for packets
addressed to a port it was listening on, and between 3.17\% and 6.67\% for an unused port, with
latency rising as the microcontroller saturated. Their conclusion was that command-and-control
over IP is unsuitable for hard real-time work on devices of this kind.

That result settles one question and opens three, and those three define this project.

\begin{enumerate}[leftmargin=1.6em,itemsep=4pt]
  \item \textbf{It measures the phenomenon but does not evaluate a mitigation.} The follow-up
        work from the same group proposes two \cite{behnke2022,behnke2023}, but neither is
        available to a developer shipping a stock ESP32: the first was evaluated against a
        simulated network interface and never built, and the second was prototyped on a Xilinx
        Zynq-7000, a considerably larger platform. Notably, \cite{behnke2023} uses the ESP32's
        second core to separate the observed process from the test harness rather than as a
        mitigation, so \emph{core pinning is nowhere evaluated as a defence}.
  \item \textbf{All the load in that line of work is applied at the IP layer.} The attack that
        actually threatens a wireless alarm is deauthentication — 802.11 management frames
        handled in the MAC layer and the radio driver, beneath the IP stack, which the device
        must process whether or not a socket is listening. Whether this interference path
        behaves like IP load is unknown.
  \item \textbf{No treatment of the isolation argument is formal, and none models an
        adversary} rather than volumetric stress.
\end{enumerate}

A second motivation emerged during the work itself and ultimately shaped the paper more than
the first. In the course of trying to verify the two standard assertions, three defects were
found \emph{in this project's own artefacts} — not in third-party code. Each had survived the
techniques one would normally rely on. That experience, rather than any single finding, is what
the research reports.

\section{Research objectives}

\begin{enumerate}[leftmargin=2.2em,itemsep=4pt,label=\textbf{O\arabic*.}]
  \item Determine whether core pinning combined with bounded queue discipline is adequate as an
        isolation mitigation on a single-chip Wi-Fi SoC, and establish precisely which part of
        that question can be settled without physical measurement.
  \item Establish, formally, whether an adversary with remote command-interface access can
        affect the safety function's control or data flow.
  \item Determine experimentally which command-path defence preserves the safety function,
        operator availability and replay resistance, treating the three as separate properties
        rather than as one.
  \item Determine where broker-side authorisation and application-layer authentication fail,
        and whether either dominates the other under credential compromise.
  \item Derive a response-time bound for the safety loop and identify which interference terms
        can be settled analytically and which require hardware.
  \item Establish whether the configuration defect found in this system generalises to
        comparable open-source projects.
  \item Specify a measurement protocol capable of falsifying the timing claim that remains
        unresolved, and publish it whether or not it can be executed.
\end{enumerate}

\section{Scope of the work}\label{sec:intro-scope}

\subsection*{What is inside the scope}

One artefact is studied: an ESP32-WROOM-32 intrusion alarm, firmware v2.2.0, whose complete
source is in the repository. The evaluation covers its command path, its safety state machine,
its queue discipline, its broker authorisation and its application-layer authentication. The
evidence is of five kinds, kept separate throughout: formal verification, host simulation,
testbed measurement against a live broker, analysis, and a static corpus study.

\subsection*{What is outside the scope, and why}

\textbf{Physical timing measurement.} No latency figure is reported anywhere in this report or
in the accompanying manuscript. The simulation harness advances a logical clock, so it provides
no information about timing whatever. The hardware protocol that would settle the question is
specified in Appendix~\ref{app:hardware} and has not been run.

\textbf{Attack tooling.} The repository publishes no procedure or tool for generating
radio-layer attacks. This is deliberate, and it limits what can be reported about the one
condition that would need such tooling.

\textbf{Generality.} One chip, one application, two framework versions, no deployment. The
corpus study in Section~\ref{sec:res-corpus} was the attempt to generalise and it returned
nothing; that outcome is reported as a null result and no prevalence claim is made from it.

\subsection*{Relationship to the accompanying manuscript}

Two documents exist and they serve different purposes. The manuscript is prepared for journal
submission: it is concise, confined to what is defensible before reviewers, and formatted to
the target journal's requirements. This report documents the whole project — including the
development environment, the implementation at module level, the experimental apparatus, and
the defects found in the project's own measurement infrastructure along the way, which a
journal paper has no room for. Where the two overlap, the numbers are identical, because both
draw on the same committed records.

\section{Main contributions}

Eight contributions were claimed in an early draft. They were reduced to five after each was
classified against the literature, because the evidence does not support eight. The reduction
is recorded in \code{docs/CONTRIBUTIONS.md}.

The central claim is this:

\begin{quote}
A security property asserted in build configuration, and a timing property asserted in design,
are both unverified claims. Each can be falsified by a method that costs far less than the
system it checks.
\end{quote}

\begin{enumerate}[leftmargin=2.2em,itemsep=5pt,label=\textbf{C\arabic*.}]
  \item \textbf{A configuration defect that costs a safety property.} A compile-time guard
        accepted TLS plus a broker login as evidence of publisher authority. The inference is
        unsound because the broker re-originates messages. Under a single leaked credential the
        alarm fired in 16.7\% of 300 trials against 100\% for the control. The guard emitted no
        diagnostic and survived review. (Chapter~\ref{ch:security})
  \item \textbf{A response-time decomposition that exposed an unbounded drain.} The
        decomposition removes three of four interference terms by construction and, in doing
        so, revealed that the safety loop had no static bound on its own execution time. The
        defect had passed 9,600 adversarial trials, 3,000,000 fuzzing iterations, exhaustive
        model checking and 129 host tests. (Chapter~\ref{ch:impl}, Section~\ref{sec:impl-drain})
  \item \textbf{A disclosed vacuous invariant, and a regression test for vacuity.} The property
        stating that unauthenticated remote disarm never takes effect was unfalsifiable in the
        specification that was shipped and checked in CI, because it assumed the refusal rather
        than deriving it. The model-checking runner now \emph{requires} a deliberately
        defective configuration to fail. (Chapter~\ref{ch:formal})
  \item \textbf{A measurement of defence independence.} A twelve-configuration ablation over
        3,600 trials shows availability, integrity and the safety function failing
        independently, and shows broker authorisation and message authentication failing under
        different compromises — reproduced against a live broker.
        (Chapter~\ref{ch:results})
  \item \textbf{A published, falsifiable measurement protocol} for the timing question that
        remains open, together with an honest negative generalisation result from a
        commit-pinned corpus of 135 projects. (Appendix~\ref{app:hardware},
        Section~\ref{sec:res-corpus})
\end{enumerate}

No new mechanism, attack, protocol or formalism is claimed. Ablation, response-time analysis
and specification mutation are all established techniques. What the work reports is that
applying them to a system's \emph{assumptions}, rather than to its running behaviour, caught
three defects that testing the running behaviour did not.

\section{Organisation of this report}

Chapter~\ref{ch:lit} reviews the relevant literature and states the research gap.
Chapter~\ref{ch:req} gives the requirements, system assumptions and threat model.
Chapter~\ref{ch:arch} presents the architecture. Chapter~\ref{ch:tools} documents the
development environment with verified versions. Chapter~\ref{ch:impl} explains the
implementation module by module with code excerpts. Chapter~\ref{ch:method} documents the
experimental methodology and Chapter~\ref{ch:results} the results and their statistical
treatment. Chapter~\ref{ch:formal} covers formal verification. Chapter~\ref{ch:security}
discusses the security evaluation. Chapter~\ref{ch:figures} catalogues the figures.
Chapter~\ref{ch:limitations} states the limitations and threats to validity, and
Chapter~\ref{ch:conclusion} concludes. Eleven appendices carry the inventories, commands and
reproduction instructions.
